\section{Source Section III.G: Wetting Dynamics}
\label{sec:wetting-dynamics}

\sourceeqrange{Eq. (3.2), figures 14--15}

\subsection*{Purpose}

This subsection connects the wetting examples back to the wall functional and
density boundary condition already derived in Section II.C and scaled in
Appendix B.

\subsection*{Reading Guide}

The source initializes a wall-attached liquid region and then compares boundary
temperature branches.  The wall density condition is the same liquid-wetting
condition used throughout Section III.  The bottom heat-flux diagnostic from
Eq. (3.18) is relevant when reading the contact-region heat absorption figure.

\subsection*{Boundary Interpretation}

The companion treats wall attraction, evaporation, condensation, and heat-flux
localization as source-stated simulation phenomena.  It does not introduce a
new contact-angle law or an independent wall-free-energy parameter beyond the
ones already recorded for Eqs. (2.33)--(2.34) and Eq. (3.2).

The useful finite objects for this scenario are boundary objects:
\[
\begin{array}{c|c|c}
\text{source anchor} & \text{object} & \text{interpretive role} \\
\hline
\text{Eq. (3.2)} & \text{scaled density wall condition}
  & \text{liquid-wetting branch inherited from II.C} \\
\text{Eq. (3.18)} & Q_b(x)
  & \text{bottom heat-flux diagnostic where }v=0 \\
\text{figures 14--15} & \text{wetting response}
  & \text{source simulation output, not derived here}
\end{array}
\]
Thus the section explains which boundary formulas to keep in view while reading
the figures.  The wall condition controls the density-gradient branch, while
\(Q_b(x)\) is a heat-flux diagnostic on the bottom wall.  These are different
boundary objects, and neither one supplies a dynamic contact-line model, an
evaporation rate law, or the unpublished pressure-tensor branch.
The finite wall-condition and bottom-heat-flux checks are covered by the
Section III mirror pair listed in Section~\ref{sec:numerical-results-overview}.
They do not reproduce wetting dynamics.


The heat-flux scale used when reading the wetting heat-absorption figure is
\begin{equation}
  Q_0=\frac{\epsilon\ell}{v_0t_0}.
  \label{eq:iii-wetting-heat-flux-normalization}
\end{equation}
The type check is
\begin{equation}
  \frac{\epsilon}{v_0}\frac{\ell}{t_0}
  =\text{energy density}\times\text{velocity},
\end{equation}
which has heat-flux units, energy per area per time.  This normalization is a
plot-reading scale; it is not a new boundary law.

\subsection*{Limitations}

No spreading law or evaporation rate is derived here.  A later numerical study
would need explicit authorization and a separate reproducibility protocol.
